\section{实验协议与评估可靠性}\label{app:exp}
\subsection{实验配置}\label{app:config}
本附录在一处列出：哪个模型扮演哪个角色、各自如何部署和解码，以及一次运行所获得的全部预算。表~\ref{tab:models} 涵盖模型，表~\ref{tab:params} 涵盖循环和评分参数。

\paragraph{评分协议。}表~\ref{tab:main} 中的每个数字都是 140 个任务上的平均值，每个任务的分数由 Qwen3.8-27B 评审对同一产物进行三次独立的回放加打分运行得到。生成器未产出项目的任务记 0 分；开发尚未完成的任务也记 0 分。基线行对冻结的 $P_0$ 打分，Play2Code 行对三十轮所交付的版本打分，\rsig{} 行则对饱和停止触发时全局质量监控器所持有的版本打分。

\begin{table}[htbp]
\caption{\textbf{按角色划分的模型。}生成器因行组而异；其余每个角色在每一行中都是同一个模型，因此行与行之间的比较是开发方法的比较，而不是其中所用模型的比较。凡是答案为关于游戏的论断（坐标、判定、评分细则条目）时，解码均采用贪心方式。部署方式对成本和可复现性都有影响——OpenRouter 上的同一个模型由十几家提供商出售，因此评审和监控器固定了各自的提供商。}
\label{tab:models}
\centering\small
\begin{tabular}{@{}llll@{}}
\toprule
角色 & 模型 & 部署方式 & 解码 \\
\midrule
生成器 & \makecell[l]{表~\ref{tab:main} 中每个\\行组各一个} & \makecell[l]{框架\\默认} & \makecell[l]{框架\\默认} \\
控制器 & GLM-5.3-Flash & OpenRouter & $T=0$，top-$p$ 1 \\
探索器 & Qwen3.8-27B & 本地 vLLM（TP $=2$） & $T=0$，top-$p$ 1 \\
验证器 & GLM-5.3-Flash & OpenRouter & $T=0$，top-$p$ 1 \\
编辑器 & GLM-5.3-Flash & OpenRouter & $T=0$，top-$p$ 1 \\
全局质量监控器 & Qwen3.8-Flash & OpenRouter & $T=0$ \\
素材生成 & Seedream-5.0-lite & OpenRouter（图像） & — \\
评估评审 & Qwen3.8-27B & 本地 vLLM & \makecell[l]{$T=0$，关闭思考，\\输出 $\leq$2048} \\
\bottomrule
\end{tabular}
\end{table}

\subsection{Token 与成本核算}\label{app:cost}
\paragraph{范围。}我们对每个任务统计两个阶段。\emph{生成}是生成器产生冻结基础 $P_0$ 的单次智能体运行。\emph{开发}是开发轮次中的每一次模型调用：决定方向的控制器、游玩游戏的探索器、检查游戏的验证器、编辑游戏的编辑器，以及全局监控器的周期性检查。所有方法都不计入评估评审的打分。

\paragraph{Token。}Tok. 一列统计计费 token：不含缓存读取的输入，加上输出。缓存读取不计入这一统计，因为长时间的智能体会话在每次调用时都会重读同一前缀；但它们仍然按缓存价格计入 Cost。

\paragraph{成本。}每次调用，
\[
\text{cost} = n_{\text{in}}\,p_{\text{in}} + n_{\text{cache}}\,p_{\text{cache}} + n_{\text{out}}\,p_{\text{out}},
\]
其中 $n_{\text{in}}$ 为未缓存的输入 token，价格采用 OpenRouter 的标价（表~\ref{tab:prices}）。开发阶段并不按此方式定价，而是汇总 API 随每次调用返回的金额——代理会为控制器、探索器、验证器、编辑器和监控器的每次调用记录这一金额；上述公式用于只保留了 token 计数的场合，即生成阶段。Cost 只涵盖开发；生成成本在表~\ref{tab:gen} 中单独报告。生成成本主要由框架而非模型决定：同一个 GPT-5.5 在 Codex 下的缓存命中率为 94\%，在 OpenGame 下为 42\%，由此产生的账单相差八倍。

\begin{table}[htbp]
\caption{\textbf{循环、监控器与评分参数。}每个参数取一个值，在本文报告的每次运行中都相同。}
\label{tab:params}
\centering\small
\begin{tabular}{@{}llr@{}}
\toprule
& 参数 & 取值 \\
\midrule
\multirow{5}{*}{轮次} & 编辑器每轮可进行的工具调用次数 & 26 \\
& 每个改进轮次的墙钟时间预算 & 1800 s \\
& 每轮放弃前的改进尝试次数 & 3 \\
& 每个声称的修复的验证采样次数 & 3 \\
& 提出编辑前的探针步数 & 8 \\
\midrule
\multirow{3}{*}{运行} & 每次运行生成的素材数 & 20 \\
& 保存的项目快照 & 每轮 \\
& 保存的回放录像 & 每 3 轮 \\
\midrule
\multirow{5}{*}{监控器} & 检查点间隔 & 每 3 轮 \\
& 每个检查点回放的演示数 & $\leq$3 \\
& 比较时每个演示读取的帧数 & 16 \\
& 候选必须超过冠军的幅度 & 0.05 \\
& 饱和停止前冠军保持不变的检查点数（$K$） & 3 \\
\midrule
\multirow{3}{*}{评分} & 每个产物的回放加打分次数 & 3 \\
& 每个演示均匀采样的帧数 & 16 \\
& 每个任务的细则条目数（$M/D/V/A$） & 12–21（中位数 18） \\
\bottomrule
\end{tabular}
\end{table}

\paragraph{采用哪个价格。}OpenRouter 上的一个模型由十几家提供商提供服务，价格相差可达两倍，而模型页面只引用其中一家的报价，对若干模型而言那恰好是全场最贵的。因此，我们按每个模型最便宜的标准端点定价（表~\ref{tab:prices}）。有两项排除使这一做法定义明确。延迟分级并不是同一服务的不同卖家，因此不予考虑：flex 分级通过延后请求使 GPT-5.5 价格减半。并且，各端点是按这些运行在其上实际会花费多少来比较的，而不仅仅比较输入价格，因为报价中缓存价格为零意味着该端点不提供提示缓存，而不是缓存读取免费——若按字面理解，一家不提供缓存的提供商就会成为一个缓存占比 42–96\% 的工作负载的最便宜卖家。

\paragraph{数据来源。}生成阶段的用量来自每次生成运行自身的记录（输入、缓存读取和输出 token），随基础游戏语料一同发布；对于 Godot 上的 Codex + GPT-5.5，只记录了语料层面的均值。开发阶段的用量由智能体与 API 之间的代理逐次调用记录。

\begin{table}[htbp]
\caption{\textbf{OpenRouter 标价（美元 / 百万 token），获取于 2026-09-16。}每个模型按其在本文运行的缓存比例下最便宜的标准端点定价；之所以注明提供商，是因为同一模型在不同提供商之间的价格相差可达两倍。}
\label{tab:prices}
\centering\small
\begin{tabular}{@{}llrrr@{}}
\toprule
模型 & 提供商 & 输入 & 缓存读取 & 输出 \\
\midrule
GPT-5.5 & OpenAI & 5.000 & 0.500 & 30.00 \\
Kimi-K2.6 & Baidu & 0.408 & 0.069 & 1.72 \\
Qwen3.8-27B & DeepInfra & 0.150 & 0.037 & 1.88 \\
GLM-5.3-Flash & DeepInfra & 0.075 & 0.015 & 0.25 \\
\bottomrule
\end{tabular}
\end{table}

\begin{table}[htbp]
\caption{\textbf{冻结基础 $P_0$ 的生成，每个任务的均值。}“计费”即基线行的 Tok.；“总计”另加缓存读取；“成本”将表~\ref{tab:prices} 应用于全部三种 token 计数。两行 GPT-5.5 是同一模型在两种框架下的结果，价格不同：Codex 会重读缓存的前缀，OpenGame 则不会。统计范围为全部 140 个任务。}
\label{tab:gen}
\centering\small
\begin{tabular}{@{}llrrrrr@{}}
\toprule
引擎 & 生成器 & 计费 & 总计 & 缓存 & 输出 & 成本 \\
\midrule
Godot & Codex + GPT-5.5 (high) & 0.26M & 3.83M & 94\% & 28k & \$\,3.77 \\
Godot & GLM-5.3-Flash & 0.44M & 11.43M & 96\% & 36k & \$\,0.20 \\
Godot & Kimi-K2.6 & 3.22M & 8.98M & 65\% & 123k & \$\,1.87 \\
Godot & Qwen3.8-27B & 6.41M & 11.62M & 46\% & 179k & \$\,1.47 \\
Godot & Qwen3.8-27B (SFT) & 0.57M & 8.21M & 95\% & 140k & – \\
Phaser & OpenGame + GPT-5.5 & 5.36M & 9.21M & 42\% & 97k & \$31.17 \\
Phaser & Qwen3.8-27B & 9.03M & 26.69M & 67\% & 292k & \$\,2.51 \\
\bottomrule
\end{tabular}
\end{table}

\subsection{回放与评分稳定性}\label{app:stability}
\paragraph{回放变异性。}GameCraft-Bench 依据观察到的游戏过程而不仅仅是源代码来为游戏打分。即使对于同一个冻结项目，由于游戏执行、输入调度和画面捕获并不完全同步，重复回放也可能产生略有不同的轨迹。这种变异的幅度取决于游戏，表~\ref{tab:family-var} 测量了它取决于哪类游戏。因此，对同一段录像反复评审并不能刻画评估过程的全部不确定性。

\paragraph{评估协议。}为了应对这种变异性，每个产物都经过三次独立的端到端运行评估。每次运行都从一次全新的回放开始，生成自己的游戏录像，并由本文所有分数都使用的同一个 Qwen3.8-27B 评审独立打分。我们报告三次运行的平均分。评审配置在所有评估中保持固定，因此重复运行主要刻画的是游戏回放引入的变异，而不是评估设置的变化。

\paragraph{观察到的稳定性。}表~\ref{tab:stability} 汇总了 529 个冻结产物上的变异，即本文中每个被打分三次的 Godot 产物：两个生成器的冻结基础、Play2Code 交付的版本，以及 \rsig{} 交付的版本。

大多数产物接近确定性。三次独立运行间极差的中位数为 0.73 分，529 个中有 225 个三次得分完全相同，四分之一的产物变化不到百分之一分。改进后的产物比冻结的产物变化略大（中位数 0.87 对 0.36），这正是一个包含更多游戏内容的构建应有的表现。

\begin{table}[htbp]
\caption{\textbf{同一冻结产物三次独立端到端评估的得分变异。}每次运行使用全新的游戏回放，并由 Qwen3.8-27B 评审重新评判。极差指三次运行中最高与最低总分之差；各列为其四分位数。}
\label{tab:stability}
\centering\small
\begin{tabular}{@{}lrrrrr@{}}
\toprule
子集 & $n$ & p25 & 中位数 & p75 & 三次相同 \\
\midrule
全部产物 & 529 & 0.00 & 0.73 & 3.33 & 225 \\
基础 $P_0$ & 266 & 0.00 & 0.36 & 2.33 & 131 \\
改进后 & 263 & 0.00 & 0.87 & 3.50 & 94 \\
\bottomrule
\end{tabular}
\end{table}

\begin{table}[htbp]
\caption{\textbf{按任务类别的回放变异}，仅列出至少有十二个已打分产物的类别，按极差中位数排序。产物集合与极差定义同表~\ref{tab:stability}。}
\label{tab:family-var}
\centering\small
\begin{tabular}{@{}lrr@{\hspace{3em}}lrr@{}}
\toprule
类别 & $n$ & 中位数 & 类别 & $n$ & 中位数 \\
\midrule
sports & 14 & 4.66 & shooter & 24 & 0.44 \\
horror & 19 & 3.37 & puzzle & 32 & 0.42 \\
openworld & 62 & 1.41 & simulation & 24 & 0.28 \\
visualnovel & 32 & 1.17 & strategy & 63 & 0.00 \\
platformer & 77 & 0.87 & rhythm & 21 & 0.00 \\
tycoon & 58 & 0.73 & idle & 16 & 0.00 \\
roguelike & 53 & 0.62 & racing & 16 & 0.00 \\
cardgame & 18 & 0.50 & & & \\
\bottomrule
\end{tabular}
\end{table}

\paragraph{哪些游戏会变化。}按任务类别拆分这 529 个产物，可以定位变异的来源，而不是把它笼统地归为“游戏动态”的属性（表~\ref{tab:family-var}）。划分并非按快慢：变化最大的三个类别是 sports、horror 和 openworld，而 racing 和 rhythm——它们与基准中任何游戏一样都是实时的——极差中位数恰好为零，strategy 亦然。变化大的类别的共同点在于，它们的评分细则答案出现得较晚。一场比赛是否获胜、一次闹鬼事件是否解决、一个世界是否被穿越，取决于二十秒的游玩恰好结束在哪里；而一个节拍是否准时落下、一个圈速计时器是否运行，则取决于回放每次都能重现的那些帧。变异追随的是某条准则的证据在会话中出现得有多晚，而不是游戏运动得有多快。

这也说明了单次运行比较在哪里最不可靠。一分的差异对节奏游戏是有意义的，对体育游戏则是噪声，这也是阅读附录~\ref{app:family} 的按类别表格时应结合该类别变异的原因之一。

\paragraph{跨评审稳健性。}为了检验主要结果是否依赖于某个特定的评估模型，我们用 GPT-5.5 重新为 40 个按类别分层抽样的 Godot 任务打分，使用与 Qwen3.8-27B 评审完全相同的回放、采样帧、任务特定细则和评分公式；只更换评审模型。

如表~\ref{tab:crossjudge} 所示，两个评审给出相同的方法层面排序：\rsig{} $>$ Play2Code $>$ 基线，且改进幅度相近。\rsig{} 相对冻结基线的增益在 Qwen3.8-27B 下为 +14.04 分，在 GPT-5.5 下为 +12.04 分；相对 Play2Code 的增益分别为 +13.47 和 +11.76；四个 bootstrap 区间都远高于零。在 40 个任务中，两位评审对“\rsig{} 优于基线”这一比较的方向在 37 个任务上一致。尽管它们的绝对分数仅中等相关（$\rho = 0.64$），但从相对比较中得出的结论在不同模型家族间保持稳定。

\begin{table}[htbp]
\caption{\textbf{40 个按类别分层的 Godot 任务上的跨评审稳健性。}Qwen3.8-27B 和 GPT-5.5 使用相同的细则为相同的回放录像打分。两位评审都保持了方法排序，并给出相近的成对改进幅度。置信区间为 95\% 任务级 bootstrap 区间。}
\label{tab:crossjudge}
\centering\small
\begin{tabular}{@{}lcc@{}}
\toprule
& \multicolumn{2}{c}{总分} \\
\cmidrule(l){2-3}
方法 & Qwen3.8-27B & GPT-5.5 \\
\midrule
基线（冻结的 $P_0$） & 49.07 & 50.34 \\
\quad + Play2Code & 49.63 & 50.62 \\
\quad + \rsig & \textbf{63.10} & \textbf{62.38} \\
\bottomrule
\end{tabular}

\vspace{6pt}
\setlength{\tabcolsep}{4pt}
\begin{tabular}{@{}lrrccc@{}}
\toprule
比较 & $\Delta$ Qwen & $\Delta$ GPT-5.5 & 95\% CI Qwen & 95\% CI GPT-5.5 & 一致数 \\
\midrule
\rsig{} $-$ 基线 & +14.04 & +12.04 & [+9.84, +18.59] & [+9.53, +14.60] & 37/40 \\
\rsig{} $-$ Play2Code & +13.47 & +11.76 & [+8.95, +18.19] & [+7.79, +16.26] & 34/40 \\
Play2Code $-$ 基线 & +0.56 & +0.28 & [$-$3.45, +4.21] & [$-$3.83, +3.83] & 24/40 \\
\bottomrule
\end{tabular}

\vspace{2pt}
{\footnotesize 一致数统计两位评审偏好同一方法的任务数。}
\end{table}

\subsection{自由游玩评估}\label{app:freeplay}
\paragraph{协议。}GameCraft-Bench 通过回放随提交项目打包的演示轨迹来评估游戏。这些轨迹是可见的开发资源，而不是隐藏的基准轨迹，\rsig{} 会复用它们在检查点之间进行受控的前后比较。尽管隐藏的评分细则、分数和评审反馈在开发过程中从不暴露，反复使用相同的演示仍可能偏向针对这些交互轨迹的特定改进。因此，我们评估这些收益在独立生成的游戏过程中是否依然存在。

我们抽取 20 个按类别分层的 Godot 任务。对每个任务，冻结基线、Play2Code 和 \rsig{} 的构建分别由盲测智能体独立游玩，每个智能体有 40 次操作和一次重置。智能体只能看到公开的游戏规格说明和正在运行的游戏——看不到提交的演示、评分细则、分数、源代码、开发历史或方法身份——并返回关于其观察的事实性报告。另一个评审对这些报告进行成对比较，共得到 60 次比较。

\paragraph{评估。}成对评审遵循主评估所用的四个质量维度和相对权重，同时只依据自由交互中观察到的证据进行评判。重要的是，回放脚本从不向游玩智能体展示：它们自行选择动作、时机和轨迹。这改变了交互分布，同时保持游戏质量的定义与主评估一致。

\paragraph{结果。}如表~\ref{tab:freeplay} 所示，\rsig{} 在 20 个任务中的 17 个上优于冻结基线、在 18 个上优于 Play2Code，而 Play2Code 相对基线并未显示出明显优势。在所有有胜负的比较中，自由游玩评估与基准排序在 59 例中的 45 例上一致。这些结果表明，\rsig{} 的收益在独立选择的游戏轨迹下依然存在，而不局限于开发和基准评估中使用的固定回放脚本。

\begin{table}[htbp]
\caption{\textbf{20 个按类别分层的 Godot 任务上的自由游玩评估。}盲测智能体独立游玩每个构建，另一个评审使用与主评估相同的质量维度和权重比较它们的报告。$p$ 为对有胜负比较所做的双侧符号检验。}
\label{tab:freeplay}
\centering\small
\begin{tabular}{@{}lccc@{}}
\toprule
比较 & 胜/负/平 & $p$ & 一致数 \\
\midrule
\rsig{} – 基线 & 17/3/0 & 0.003 & \multirow{3}{*}{45/59} \\
\rsig{} – Play2Code & 18/2/0 & $<$ 0.001 & \\
Play2Code – 基线 & 9/11/0 & 0.824 & \\
\bottomrule
\end{tabular}

\vspace{2pt}
\parbox{0.85\linewidth}{\footnotesize 胜/负/平为第一个方法的胜、负、平局数。一致数统计自由游玩评估与基准得分偏好同一构建的有胜负比较数。}
\end{table}

\subsection{主要结果的统计可靠性}\label{app:stat}
\paragraph{配对的任务级不确定性。}主表中的所有比较都按任务配对。我们对 140 个基准任务进行 20,000 次 bootstrap 重采样，以估计总分平均差 $\Delta$ 的不确定性，并额外报告对配对任务级差异的 Wilcoxon 符号秩检验。我们还单独对每个固定任务的三次回放加打分运行进行重采样，以衡量评估噪声；这些区间远窄于任务 bootstrap 区间，说明基准任务之间的差异主导了回放和评审噪声。

\paragraph{主要比较。}在表~\ref{tab:sig} 包含的每种引擎–生成器设置下，\rsig{} 都显著优于其冻结的起点，增益为 +8.8 到 +14.3 个总分点，任务 bootstrap 区间远高于零。与 Play2Code 相比，改进在两个 Godot 生成器以及 Phaser 上的 GPT-5.5 中都是显著的；Phaser 上 Qwen3.8-27B 的比较仍不能确定（+1.55，$p = 0.394$）。

\begin{table}[htbp]
\caption{\textbf{主要结果的配对显著性。}$\Delta$ 为任务级总分的平均差。置信区间来自对基准任务进行的 20,000 次配对 bootstrap 重采样；$p$ 来自 Wilcoxon 符号秩检验。胜/负统计第一个方法得分更高/更低的任务数。}
\label{tab:sig}
\centering\small
\begin{tabular}{@{}lrccc@{}}
\toprule
比较 & $\Delta$ & 95\% CI & $p$ & 胜/负 \\
\midrule
\multicolumn{5}{@{}l}{\textit{Godot}} \\
GPT-5.5：\rsig{} – 基线 & +14.26 & [11.49, 17.06] & $<10^{-4}$ & 112/16 \\
GPT-5.5：\rsig{} – Play2Code & +13.79 & [10.89, 16.65] & $<10^{-4}$ & 118/19 \\
Qwen3.8-27B：\rsig{} – 基线 & +10.70 & [7.75, 13.77] & $<10^{-4}$ & 100/31 \\
Qwen3.8-27B：\rsig{} – Play2Code & +7.24 & [4.49, 10.04] & $<10^{-4}$ & 92/41 \\
\midrule
\multicolumn{5}{@{}l}{\textit{Phaser}} \\
GPT-5.5：\rsig{} – 基线 & +8.77 & [6.13, 11.51] & $<10^{-4}$ & 101/35 \\
GPT-5.5：\rsig{} – Play2Code & +3.11 & [0.53, 5.61] & 0.008 & 83/53 \\
Qwen3.8-27B：\rsig{} – 基线 & +10.21 & [6.99, 13.70] & $<10^{-4}$ & 95/36 \\
Qwen3.8-27B：\rsig{} – Play2Code & +1.55 & [$-$1.50, 4.66] & 0.394 & 73/59 \\
\bottomrule
\end{tabular}
\end{table}

\subsection{自适应方向策略评估}\label{app:adaptive}
\paragraph{实验设置。}我们在 16 对从相同冻结基础项目初始化的游戏上，比较自适应开发与轮转方向策略。每种策略运行 30 个开发轮次，使用相同的智能体模型、工具和每轮预算；唯一的区别在于如何选择开发重点。图~\ref{fig:adaptive}a,b 使用在开发过程中对两种策略都不可见的信号来分析所得轨迹。我们将“构建失败”定义为所得项目无法产出可运行游戏的轮次；随后的轮次相对于该失败事件对齐到偏移 +1。对于按状态条件的分析，我们另外使用一份保留的诊断细则来识别当前落后的质量维度。该细则仅用于事后分析，从不提供给开发循环。图~\ref{fig:adaptive}a 中的阴影区域为相应失败事件或分组开发轮次上的 $\pm1$ 标准误。

\paragraph{盲测成对评估。}为比较两种策略产出的最终游戏，我们使用它们第 30 轮的检查点进行了 63 次盲测成对比较。每次比较中，同一条演示轨迹在两个构建上回放，所得观察并排呈现给评估者。左右位置在每次比较中独立随机分配，且不透露任何关于底层方向策略的信息。评估者报告总体偏好，以及在机制、深度、视觉和美术各维度上的偏好；允许平局。为检查位置偏差，一部分比较在交换左右后被再次评估。这些重复判断仅用作一致性检查，不计入报告的 63 次比较。成对评估者独立于 GameCraft-Bench 基准评估器，无法访问其分数、细则或反馈。

\subsubsection{成对评判准则}
以下完整给出交给第~5 节盲测成对评审的准则。它们刻意不同于基准的评分细则：评审依据画面所示比较两个构建，而基准则依据任务为单个构建打分。规则 1 起的作用最大——没有它，一个坏掉的构建也能凭它确实渲染出来的那些帧拿到分数。

\begin{promptbox}{pbg5}{比较细则.\ 盲测成对评审}
\begin{PV}
你是一名资深游戏评测者。你要比较同一个游戏的两个构建，
游戏 1 和游戏 2。二者都由同一个起始游戏和同一份设计文档
开发而来，设计文档见下文。每个构建都用完全相同的脚本化
输入（一个“演示”）游玩过，你看到的是从该录像中以相同
固定间隔、按时间顺序采样的画面。只依据画面中可见的内容
进行评判。

对下面的每个维度，判断哪个构建更好："1"、"2" 或 "tie"。

## 核心机制
脚本化输入是否明显产生了设计文档所描述的游戏玩法？
- 更好：操作在屏幕上有清晰的效果（移动、攻击、放置、
  状态变化、分数或资源变化）；能观察到设计文档的核心循环。
- 更差：输入没有可见效果；游戏停在标题或菜单画面；
  场景加载失败；屏幕空白或冻结。

## 内容深度
在这个演示中，设计好的游戏有多少内容出现并被触及？
- 更好：设计文档中更多不同的机制、敌人、关卡、事件、
  进程或胜负状态出现。
- 更差：单一的重复画面；占位内容；UI 中宣布了却从未
  展示的系统。

## 功能性视觉
玩家能否读懂游戏状态？
- 更好：HUD、反馈和重要物体清晰可读，并与背景明显区分；
  文字没有被截断或重叠；精灵和动画帧与帧之间稳定。
- 更差：文字无法阅读或相互重叠；精灵闪烁或不一致；关键
  物体被其他物体遮挡；重要事件没有可见反馈。

## 呈现与美术
它看起来像一个连贯、完成的游戏吗？
- 更好：各画面风格一致；构图有意为之；素材与设计文档的
  主题相符。
- 更差：该有美术的地方是基本图形；风格混杂或冲突；视觉
  噪声；让游戏更难读懂的装饰。

## 总体
综合考虑，这份设计文档的玩家会更喜欢哪个构建？

## 规则
1. 在这些画面中坏掉的构建（空白、崩溃、冻结、卡在演示
   本应离开的画面上），在它未能展示的每个维度上都判负，
   无论它其他画面看起来多好。
2. 更多素材、更多特效或更繁忙的画面本身并不更好。只有
   当它们服务于设计并保持游戏可读时才予以偏好。
3. 忽略构建显示在哪一侧。游戏 1 和游戏 2 的顺序是随机的。
4. 当两个构建在某个维度上没有明显差异，或二者同样失败时，
   回答 "tie"。
5. 每条理由都必须指向可见的内容，点明构建和帧（例如
   “游戏 2，第 7 帧：分数计数器与计时器重叠”）。
\end{PV}
\end{promptbox}

\subsection{智能体验证审计}\label{app:audit}
\paragraph{审计协议。}我们在由 GPT-5.5 初始化的 40 个游戏的开发轨迹上评估验证的可靠性。我们从这些轨迹中抽取改进前和改进后的决策，并将与每个决策相关的完整交互证据交给 Claude Opus 5\citep{anthropic2026}进行独立裁定。每个案例评估两次；两次判断之间的分歧由人工复核以得到最终标签。该评估者仅用于验证审计，与 GameCraft-Bench 评估器及其分数、细则和反馈完全隔离。

\paragraph{改进前溯据。}我们评估所提出的改进目标是否得到编辑前可用证据的支持。对每个被审计的会话，若观察到的游戏证据支持所声称的问题或改进机会，则该目标被标记为\emph{有据}；若缺乏支持或与证据矛盾，则标记为\emph{无据}。我们报告有据精确率（可判定目标中有据者的比例）以及每轮无据目标的平均数。我们将 \rsig{} 使用的智能体证据溯据流程与一个在相同开发上下文中运行的自由形式评论者进行比较。

\paragraph{改进后失败检测。}我们进一步评估在执行更新后的游戏之后，验证能否检测出不成功的编辑。当预期的变化在游戏中无法观察到，或该编辑引入了退化时，一次改进被标记为\emph{不成功}。以不成功的改进为正类，我们在 48 个经裁定的轮次上报告失败召回率、特异度和平衡准确率。我们将基于回放的智能体验证与一个仅检查构建的基线进行比较，后者只检查更新后的项目能否成功执行，而不检查其交互行为。

\subsection{多阶段指引评估}\label{app:multistage}
第~5 节检验的是开发的单个阶段。本附录说明我们如何检验第二和第三阶段是否持续带来回报：每对阶段冠军都通过盲测成对游玩进行比较，因此较晚的阶段必须作为一个游戏被偏好，而不仅仅是得分更高。

\begin{wrapfigure}{r}{0.46\linewidth}
\vspace{-10pt}
\centering
\includegraphics[page=11,width=\linewidth]{\figpdf}
\caption{每对中较晚阶段的盲测成对胜率，附 $\pm1$ 标准误。}
\label{fig:multistage}
\vspace{-8pt}
\end{wrapfigure}

\paragraph{比较什么。}三对构建，每个都是全局质量监控器在某阶段结束时交付的冠军：阶段 2 对阶段 1、阶段 3 对阶段 2，以及阶段 3 对阶段 1。一对中的两个构建来自同一个游戏，因此每次比较都在任务内部进行，双方使用相同的生成器、相同的任务规格说明和相同的每阶段预算。结果显示，较晚的阶段始终比较早的阶段更受偏好（图~\ref{fig:multistage}），表明高层指引能够在饱和之后反复重新打开开发的提升空间，并支持持续的多阶段演化。

\paragraph{为何采用成对比较。}一个修复了坏掉机制的阶段，和一个添加了无用菜单的阶段，可能使基准得分变动相同的幅度。因此，本评估问的是玩家更愿意拥有哪个构建，而不使用基准得分作为结果。

\paragraph{三个智能体，各看一个构建。}一次比较使用三个智能体。两个游玩智能体各获得一个构建，从不看到另一个，也看不到任何分数、源文件、轮次编号或开发历史；每个智能体返回一份固定格式的事实报告——游戏能否启动、核心循环能否完成、输入如何响应、渲染出了什么、什么缺失或坏掉——并被明确禁止给出评判。第三个评审智能体只阅读这两份报告（顺序在每次比较中随机化并记录），返回“报告 1 更好”“报告 2 更好”或“平局”，并附上置信等级。评审被要求将一个具体点名的缺陷看得比笼统印象更重，并将关于评审框架的任何内容——一次调用推进多少游戏时间、一个构建启动要多久、操作预算——视为两个构建共有、因而不构成对任何一方的证据。

\paragraph{预算相等。}每个游玩智能体都在 40 次操作和一次重置的固定预算下，通过附录~\ref{app:director} 的评审工具驱动其构建，其中 \texttt{wait}、\texttt{look} 和 \texttt{state} 免费。会话以评审者 id 为键，重新打开的会话会带着已花费的预算继续，因此每个智能体都获得该次比较独有的 id，并在打开后立即检查自己的操作计数器；报告计数非零的会话会被丢弃，该比较重新运行。在不相等预算下比较两个构建，得到的结论追随的是预算而不是构建，这就是为什么这项检查是协议的一部分而非事后补充。

\paragraph{报告方式。}当评审偏好属于较晚阶段的那份报告时，该比较计为较晚阶段胜出。平局单独报告，并从胜率中排除而不是对半拆分，因为平局表明两份报告无法区分这两个构建。胜率附带 Wilson 形式的 $\pm1$ 标准误，可在接近端点时使区间保持在 $[0,1]$ 之内；每对只有几十个游戏，这些区间较宽，正文也如实说明这一点，而不是从相互重叠的柱子中读出排名。
